% Plantilla del tono institucional: portada sencilla, índices, tablas y figuras numeradas.
% Reemplaza todo lo que está entre <...>. Compila con: python <skill>/build.py <este archivo>
\documentclass[11pt,a4paper]{article}

\usepackage[spanish,es-noshorthands,es-tabla]{babel}
\usepackage{fontspec}
\setmainfont{texgyrepagella}[Extension=.otf, UprightFont=*-regular, BoldFont=*-bold,
  ItalicFont=*-italic, BoldItalicFont=*-bolditalic]
\setsansfont{texgyreheros}[Extension=.otf, UprightFont=*-regular, BoldFont=*-bold,
  ItalicFont=*-italic, BoldItalicFont=*-bolditalic]
\usepackage[hmargin=2.4cm, top=2.6cm, bottom=2.4cm, headheight=14pt]{geometry}
\usepackage{microtype}
\usepackage{xcolor}
\usepackage{graphicx}
\usepackage{tikz}
\usetikzlibrary{positioning,arrows.meta,fit,backgrounds,calc,shapes.geometric}
\usepackage{booktabs}
\usepackage{tabularx}
\usepackage{longtable}
\usepackage{array}
\usepackage{needspace}
\usepackage{float}
\newcolumntype{Y}{>{\raggedright\arraybackslash}X}
% Primera columna en negrita y en bandera, sin partir palabras.
\newcolumntype{B}[1]{>{\raggedright\arraybackslash\bfseries\hyphenpenalty=10000}p{#1}}
\usepackage{enumitem}
\usepackage{titlesec}
\usepackage{caption}
\usepackage{fancyhdr}
\usepackage{xurl}
\usepackage[hidelinks]{hyperref}
\usepackage{tocloft}

\definecolor{tinta}{HTML}{111111}
\definecolor{marino}{HTML}{000000}
\definecolor{gris}{HTML}{6E6E6E}
\definecolor{linea}{HTML}{BDBDBD}
\definecolor{fondo}{HTML}{F2F2F2}
\definecolor{fondo2}{HTML}{E4E4E4}
\color{tinta}

\titleformat{\section}{\Large\bfseries\color{marino}}{\thesection.}{0.6em}{}
\titleformat{\subsection}{\large\bfseries}{\thesubsection.}{0.5em}{}
\titleformat{\subsubsection}{\normalsize\bfseries}{\thesubsubsection.}{0.5em}{}
\titlespacing*{\section}{0pt}{18pt}{6pt}
\titlespacing*{\subsection}{0pt}{12pt}{4pt}
\setcounter{secnumdepth}{3}
\setcounter{tocdepth}{2}
\setlist{itemsep=3pt,topsep=4pt,leftmargin=1.4em}
\setlength{\parindent}{0pt}
\setlength{\parskip}{6pt}
\DeclareCaptionFont{gris}{\color{gris}}
\captionsetup{font={small,gris}, labelfont={bf,color=tinta}, labelsep=period}
\captionsetup[table]{position=top, skip=4pt}
\renewcommand{\arraystretch}{1.25}
\renewcommand{\cftsecleader}{\cftdotfill{\cftdotsep}}
\renewcommand{\cfttoctitlefont}{\Large\bfseries}
\renewcommand{\cftloftitlefont}{\Large\bfseries}
\renewcommand{\cftlottitlefont}{\Large\bfseries}

\pagestyle{fancy}
\fancyhf{}
\fancyhead[L]{\small\color{gris}<Institución · dependencia>}
\fancyhead[R]{\small\color{gris}<Título corto>}
\fancyfoot[C]{\small\color{gris}\thepage}
\renewcommand{\headrulewidth}{0.4pt}
\renewcommand{\headrule}{\color{linea}\hrule height 0.4pt\vspace{0pt}}
\fancypagestyle{plain}{\fancyhf{}\renewcommand{\headrule}{}}

\tikzset{
  caja/.style={draw=linea, fill=white, rounded corners=2pt, align=center,
    execute at begin node={\hyphenpenalty=10000\exhyphenpenalty=10000},
    font=\sffamily\small, inner sep=5pt, minimum height=1.0cm, text=tinta},
  clave/.style={caja, draw=marino, line width=0.8pt},
  persona/.style={caja, fill=fondo2, draw=marino, rounded corners=8pt},
  externo/.style={caja, fill=fondo, dashed},
  datos/.style={cylinder, shape border rotate=90, aspect=0.18, draw=marino, fill=white,
    align=center, font=\sffamily\small, minimum height=1.2cm, minimum width=2.4cm, inner sep=3pt},
  flecha/.style={-{Stealth[length=2.4mm]}, line width=0.7pt, draw=gris},
  flechad/.style={flecha, dashed},
  etq/.style={font=\sffamily\scriptsize, text=gris, align=center, fill=white, inner sep=1pt},
  limite/.style={draw=linea, dashed, rounded corners=4pt, inner sep=8pt},
  nota/.style={font=\sffamily\footnotesize\itshape, text=gris, align=center},
  estado/.style={draw=marino, fill=white, rounded corners=7pt, font=\sffamily\small,
    inner sep=4pt, minimum height=0.8cm, minimum width=2.1cm, align=center},
  entidad/.style={draw=marino, fill=white, font=\sffamily\scriptsize, align=left,
    inner sep=3pt, text width=#1},
  entidad/.default=3.1cm,
}

\begin{document}

% Portada: escudo, sede, tipo de documento, título, subtítulo, autores, ciudad y mes.
\begin{titlepage}
\thispagestyle{empty}
\includegraphics[width=6.2cm]{escudo-unal.pdf}\par
\vspace{4pt}
{\sffamily\small\color{gris} <Sede>\par}
\vspace{5.2cm}
{\sffamily\footnotesize\bfseries\color{gris} <TIPO DE DOCUMENTO>\par}
\vspace{4pt}
{\color{marino}\rule{2.2cm}{1.6pt}\par}
\vspace{14pt}
{\fontsize{30}{36}\selectfont\bfseries\raggedright\hyphenpenalty=10000 <Título,\ partido a mano si es largo>\par}
\vspace{14pt}
{\large\color{gris} <Subtítulo: de qué trata el documento en una o dos líneas>\par}
\vfill
{\color{linea}\rule{\linewidth}{0.4pt}\par}
\vspace{8pt}
{\sffamily\small\begin{tabular}{@{}l@{}}
\textbf{<Autores>}\
\color{gris} <Rol de los autores>\[6pt]
\color{gris} <Ciudad>, <mes de año>
\end{tabular}\par}
\end{titlepage}

\tableofcontents
\clearpage
\listoffigures
\vspace{1.2cm}
\listoftables
\clearpage

\section{Introducción}

\subsection{Objeto}
<Qué describe el documento y para qué decisión o trámite sirve.>

\subsection{Alcance}
<Qué cubre y qué queda por fuera.>

\subsection{Fuentes y fecha de corte}
<De dónde salen las cifras. Todas las mediciones tienen \textbf{corte en <mes de año>}. Dónde está el
código o la documentación técnica, si el lector la necesita.>

\section{<Sección>}

\begin{table}[H]
\caption{<Qué muestra la tabla, en una oración>.}
\small
\begin{tabularx}{\linewidth}{@{}B{3.4cm}Y@{}}
\toprule
\normalfont\bfseries <Columna> & \textbf{<Columna>} \\
\midrule
<Fila> & <Texto> \\
\bottomrule
\end{tabularx}
\end{table}

\begin{figure}[H]
\centering
\resizebox{\linewidth}{!}{%
\begin{tikzpicture}[node distance=8mm and 12mm]
  \node[caja] (a) {<Componente>\\\footnotesize\color{gris} <detalle>};
  \node[clave, right=of a] (b) {\textbf{<Componente clave>}};
  \draw[flecha] (a) -- (b);
\end{tikzpicture}}
\caption{<Nombre de la vista>. <Qué debe ver el lector en el diagrama>.}
\end{figure}

\section*{Referencias}
\addcontentsline{toc}{section}{Referencias}
\small
\begin{itemize}[leftmargin=1.2em]
  \item <Autor o entidad (año). \emph{Título}. \url{https://...}>
\end{itemize}

\end{document}
